\bigskip
\raggedright
\textbf{\MakeUppercase{Resumo}}

Os Modelos de Linguagem de Grande Escala (Large Language Models – LLMs) revolucionaram a forma como interagimos com a inteligência artificial. Mas, junto com essas vantagens, também surgem os desafios. Assim como qualquer tecnologia poderosa, os LLMs também podem ser explorados de forma indevida; é aí que entra o Jailbreak (conjunto de técnicas buscando burlar as restrições de segurança dos LLMs, fazendo com que respondam a comandos que normalmente seriam bloqueados). O objetivo deste artigo é analisar como o Jailbreak evoluiu de otimizações matemáticas para técnicas de furtividade semântica, persuasão humana, manipulação estrutural e pré-preenchimento de respostas via API. O trabalho é estruturado com base em uma abordagem exploratória fundamentada em levantamento bibliográfico. Inicialmente, apresenta-se a evolução para a furtividade semântica e o uso da persuasão psicológica (seções 2 e 3) Em seguida, abordam-se a manipulação da estrutura da linguagem, o controle por pré-preenchimento de respostas e os procedimentos metodológicos (seções 4, 5 e 6). Por fim, apresentam-se as conclusões do estudo. Os resultados apontam que a combinação de técnicas de pré-preenchimento com persuasão atinge taxas de sucesso de até 88\% a 89 \%,  evidenciando a fragilidade dos alinhamentos atuais e a necessidade de defesas baseadas em IA.

\palavraschave{Modelos de Linguagem de Grande Escala (LLMs). Jailbreak. Segurança em IA.}
